% !PRINTED 507
\weberpage{507}
6. \gesperrt{Ist $\omega$ ein Functional in $\Omega$ und giebt es $m$ Grössen $\omega_1, \omega_2, \ldots, \omega_m$, die nicht alle verschwinden, von der Art, dass die $m$ Producte $\omega\omega_i$ für $i=1,2,\ldots,m$ in die Form gesetzt werden können:}
\weq{5}{\omega\,\omega_i = A_{1,i}\omega_1 + A_{2,i}\omega_2 + \cdots + A_{m,i}\omega_m,}
\gesperrt{worin die $m^2$ Symbole $A_{k,i}$ ganze rationale Functionale sind, so ist $\omega$ ein ganzes Functional.}

Hierin ist $m$ irgend eine natürliche ganze Zahl. Die $\omega_i$ sind in der Anwendung immer Zahlen oder Functionale in $\Omega$, jedoch ist für die Gültigkeit des Satzes nur die Ausführbarkeit der in (5) angedeuteten Multiplication wesentlich.

Der Satz ist eine einfache Folge der Definition der ganzen Functionale; denn da die $\omega_i$ nicht alle verschwinden, so muss nach dem Determinantensatze (Bd. I, §.\,24, II.)
\weq{6}{\begin{vmatrix}
A_{1,1}-\omega, & A_{2,1}, & \ldots, & A_{m,1} \\
A_{1,2}, & A_{2,2}-\omega, & \ldots, & A_{m,2} \\
\hdotsfor{4} \\
A_{1,m}, & A_{2,m}, & \ldots, & A_{m,m}-\omega
\end{vmatrix} = 0}
sein. Durch Ordnen nach Potenzen von $\omega$ ergiebt sich hieraus eine Gleichung
\weq{7}{\omega^m + A_1\omega^{m-1} + \cdots + A_m = 0,}
worin die Coëfficienten $A_1, A_2, \ldots$ durch Addition und Multiplication aus den $A_{i,k}$ zusammengesetzt sind und daher nach 1. ganze rationale Functionale sind; demnach ist auch $\omega$ ein ganzes Functional.

Für den letzten Coëfficienten $A_m$ in der Gleichung (7) erhalten wir den Ausdruck
\weq{8}{(-1)^mA_m = \Sigma \pm A_{1,1}A_{2,2}\ldots A_{m,m},}
\gesperrt{der also gleich dem Producte der sämmtlichen Wurzeln der Gleichung (7) ist.}

Der Satz 4. ist wichtig als Kennzeichen für ganze Functionale; er dient uns hier zum Beweise des folgenden Satzes: % [sic: Weber prints Satz 4.; the determinant criterion just proved is Satz 6.]

7. \gesperrt{Ist $\Psi(x,y,\ldots)$ eine ganze Function, deren Coëfficienten ganze rationale Zahlen oder Functionale, aber frei von den Variablen $x,y,\ldots$ sind, sind ferner $\alpha,\beta,\ldots$ ganze Functionale in $\Omega$, so ist}
\weq{9}{\omega = \Psi(\alpha,\beta,\ldots)}
\gesperrt{auch ein ganzes Functional.}
